The space group $I\, 2_1\, 2_1\, 2_1$ (ITA No.~24) is generated by:
$$
G_{24} = \left\langle\left(\begin{array}{rrr|r}
-1 & 0 & 0 & 1 \\
0 & -1 & 0 & -\frac{1}{2} \\
1 & 0 & 1 & 0 \\
\end{array}\right) , \left(\begin{array}{rrr|r}
-1 & 0 & 0 & 0 \\
1 & 1 & 0 & \frac{1}{2} \\
0 & 0 & -1 & \frac{1}{2} \\
\end{array}\right) , e_1 , e_2 , e_3\right\rangle.
$$
% In this case, the matrices $\mathcal{X}$ satisfying (\textit{PC}) have a wide range of possible forms:
% $$\left(\begin{array}{rrr}
% x_{0} & 2 \, x_{0} & 0 \\
% x_{1} & -x_{0} & x_{0} + 2 \, x_{1} \\
% x_{2} & -x_{0} & 0
% \end{array}\right) , \left(\begin{array}{rrr}
% -2 \, x_{0} + x_{1} & 0 & 0 \\
% x_{2} & -2 \, x_{0} + x_{1} + 2 \, x_{2} & 0 \\
% x_{0} & 0 & x_{1}
% \end{array}\right) , \left(\begin{array}{rrr}
% x_{0} & 0 & 0 \\
% -\frac{1}{2} \, x_{0} + \frac{1}{2} \, x_{1} & 0 & x_{1} \\
% -\frac{1}{2} \, x_{0} + \frac{1}{2} \, x_{2} & x_{2} & 0
% \end{array}\right)$$
% $$\left(\begin{array}{rrr}
% -x_{0} & -2 \, x_{0} & 0 \\
% x_{1} & x_{0} & 0 \\
% \frac{1}{2} \, x_{0} + \frac{1}{2} \, x_{2} & x_{0} & x_{2}
% \end{array}\right) , \left(\begin{array}{rrr}
% x_{0} & 0 & 2 \, x_{0} \\
% x_{1} & x_{0} + 2 \, x_{1} & -x_{0} \\
% x_{2} & 0 & -x_{0}
% \end{array}\right) , \left(\begin{array}{rrr}
% x_{0} - 2 \, x_{1} & 0 & 2 \, x_{0} - 4 \, x_{1} \\
% x_{2} & 0 & -x_{0} + 2 \, x_{1} \\
% x_{1} & x_{0} & -x_{0} + 2 \, x_{1}
% \end{array}\right)$$

% For the sake of succinctness, we consider only two of them. Firstly, take a closer look on the matrix
% $A_1 = \left(\begin{array}{rrr}
% x_{0} & 2 \, x_{0} & 0 \\
% x_{1} & -x_{0} & x_{0} + 2 \, x_{1} \\
% x_{2} & -x_{0} & 0
% \end{array}\right).$
% We omit the full system of congruences (\ref{eqn:cocycle_condition}) due to its complexity. Instead, here is the
% smaller system (\ref{eqn:small_cong_system}) over integral variables $x_i$:
% $$\left(\begin{array}{r}
% x_{0} + 2 \, x_{2} \\
% -\frac{1}{2} \, x_{0} + x_{1} - x_{2} \\
% -\frac{3}{2} \, x_{0} - x_{1} - x_{2} \\
% 2 \, x_{0} + 2 \, x_{2} \\
% -x_{0}
% \end{array}\right) \equiv \left(\begin{array}{r}
% 0 \\
% \frac{1}{2} \\
% \frac{1}{2} \\
% -2 \\
% 1
% \end{array}\right) \quad \mod 1$$
% The solution is the set of matrices, satisfying the (\textit{PC})-(\textit{CC}) conditions:
% $$
% A_1'=\left(\begin{array}{rrr}
% -2 \, y_{0} - 2 \, y_{1} - 2 \, y_{2} - 1 & -4 \, y_{0} - 4 \, y_{1} - 4 \, y_{2} - 2 & 0 \\
% y_{0} + y_{1} & 2 \, y_{0} + 2 \, y_{1} + 2 \, y_{2} + 1 & -2 \, y_{2} - 1 \\
% 2 \, y_{0} + y_{1} + y_{2} & 2 \, y_{0} + 2 \, y_{1} + 2 \, y_{2} + 1 & 0
% \end{array}\right)
% $$
% The $A_1'$ has a reduced characteristic polynomial since $\tau(A_1') = \tr(A_1') = 0$, and SRD has the form (\ref{eqn:space_case2}).
% The determinant $\det(A_1')$ is a qubic polynomial of three variables No. 7 from the Table \ref{tab:case2_srdegrees0}:

For the sake of succinctness, we consider only two subsets of $SRD$ for $G_{24}$. The first is characterised by a
qubic polynomial of three variables No.~7 from the Table~\ref{tab:case2_srdegrees0}:
\begin{align*}
f_\sigma(y_0, y_1, y_2) &= 4 \, {\left(2 \, y_{0} - 1\right)} y_{2}^{2} + 4 \, y_{0}^{2} + 2 \, {\left(2 \, y_{0} - 1\right)} y_{1} \\
& + 4 \, {\left(2 \, y_{0}^{2} + {\left(2 \, y_{0} - 1\right)} y_{1} + y_{0} - 1\right)} y_{2} - 1.
\end{align*}
We can factorize it:
$$
f_\sigma(y_0, y_1, y_2) = (2y_0 - 1)(2y_2 + 1)(2y_0 + 2y_1 + 2y_2 + 1).
$$
Now it is clear that $f_\sigma$ can attain any odd integer $n$ taking $y_0 = \frac{n+1}{2}, \, y_1 = -y_0, \, y_2 = 0$.
As a result, the one of subsets of $SRD$, that coincides with a subset of $SCD$:
$$
SCD_1 = SRD_1 = \{ 2k + 1 \, : \, k \in \mathbb{N}\}.
$$

% Now consider the matrix $A_2 = \left(\begin{array}{rrr}
% \frac{1}{2} \, x_{0} & 0 & x_{0} \\
% x_{1} & \frac{1}{2} \, x_{0} + 2 \, x_{1} & -\frac{1}{2} \, x_{0} \\
% x_{2} & 0 & -\frac{1}{2} \, x_{0}
% \end{array}\right) $. It contains rational coefficients, so we have to rewrite it to satisfy (\textit{LC}):
% $$A'_2 = \left(\begin{array}{rrr}
% z_{0} & 0 & 2 \, z_{0} \\
% z_{1} & z_{0} + 2 \, z_{1} & -z_{0} \\
% z_{2} & 0 & -z_{0}
% \end{array}\right).$$
% Now the congruence system (\ref{eqn:small_cong_system}) over integral variables $z_i$ is similar to the previous one:
% $$\left(\begin{array}{r}
% 0 \\
% -\frac{1}{2} \, z_{0} + z_{1} - z_{2} \\
% -\frac{3}{2} \, z_{0} - z_{1} - z_{2} \\
% -2 \, z_{0} - 2 \, z_{2} \\
% -z_{0}
% \end{array}\right) \equiv \left(\begin{array}{r}
% 0 \\
% \frac{1}{2} \\
% \frac{1}{2} \\
% 2 \\
% 1 
% \end{array}\right) \quad \mod 1,$$
% which yields the matrix satisfying all the (\textit{PC})-(\textit{CC}) conditions:
% $$A_2'' = \left(\begin{array}{rrr}
% -2 \, y_{0} - 1 & 0 & -4 \, y_{0} - 2 \\
% y_{0} - y_{1} & -2 \, y_{1} - 1 & 2 \, y_{0} + 1 \\
% 2 \, y_{0} - y_{1} - y_{2} & 0 & 2 \, y_{0} + 1
% \end{array}\right).$$

% In contrast to $A_1'$, $A_2''$ is block-triagonal up to the renaming of the basis vectors.
% Although the minor $C'$ for $A_2'$ has only two variables and appears in the plane group No. 9
% (check Table~\ref{tab:planar_srdegrees}), the respective minor $C''$ for $A_2''$ suddenly has three variables. The $SRD$
% is described by the following polynomials:

The second subset is characterised by a triple of polynomials that can be found in the Table~\ref{tab:case1_srdegrees0}:
$$
p_\sigma(y_0, y_1, y_2) = -2 y_1 - 1,
$$
$$
q_\sigma(y_0, y_1, y_2) = 4 y_0^2 - 2(2 y_0 + 1) y_1 - 2 (2 y_0 + 1) y_2 - 1,
\quad g_\sigma = 0.
$$
% This triple can be found in the Table \ref{tab:case1_srdegrees0}.
Since $g_\sigma = 0$, the condition on eigenvalues by the Proposition \ref{propos:planar-srdegrees} reduces
to $q_\sigma \neq -1$. We can also factorize $q_\sigma$ to inspect what values $(p_\sigma*q_\sigma)$ can attain in
(\ref{eqn:space_case2}):
$$
q_\sigma = (2y_0 + 1)(2y_0 - 2y_1 - 2y_2 - 1)
$$
Similarly to the first case, $\det(A) = p_\sigma*q_\sigma$ attains every odd integer $n > 1$ via
$y_1 = -\frac{n+1}{2}, y_2 = \frac{n-1}{2}, y_0 = 0$. It does not violate the condition on eigenvalues, since 
$q = 1 \neq -1$.
Moreover, we can now observe that the substitution $y_0 \mapsto -y_1, \, y_1 \mapsto -y_2 - 1, \, y_2 \mapsto y_0$ into the
polynomial $f_\sigma$ yields exactly the $p_\sigma*q_\sigma$. As a result:
$$
SCD_2 = SRD_2 = \{ 2k + 1 \, : \, k \in \mathbb{N}\}.
$$

Note that these two subset do not represent the entire sets $SCD$ and $SRD$. To compute them completely, one must
repeat the provided procedure for every triple in the Tables~\ref{tab:case1_srdegrees0}-\ref{tab:case1_srdegrees286}
and every qubic polynomial in the Table~\ref{tab:case2_srdegrees0}.
